% !PRINTED 452
\weberpage{452}
Das ganze System geht durch 168 Substitutionen der Gruppe in sich selbst über.

\[
\begin{aligned}
&\text{Auf jeder Linie } A \text{ liegen vier Punkte } a. \\
&\text{Durch jeden Punkt } a \text{ gehen vier Linien } A. \\
&\text{Auf jeder Linie } T \text{ liegen drei Punkte } a. \\
&\text{Durch jeden Punkt } t \text{ gehen drei Linien } A.
\end{aligned}
\]

\gesperrt{Die Punkte} $a$ \gesperrt{und} $t$ \gesperrt{bilden zusammen das vollständige System aller Schnittpunkte der Linien} $A$\gesperrt{.}

Denn 21 gerade Linien schneiden sich in 210 Punkten, und von diesen fallen drei oder sechs zusammen, wenn drei oder vier dieser Linien durch einen Punkt gehen. Es ist aber
\[
210 = 21.6 + 28.3.
\]

\gesperrt{Auf jeder Linie} $A$ \gesperrt{liegen vier Punkte} $t$\gesperrt{.}

Denn da durch jeden Punkt $t$ drei Linien $A$ gehen, und auf jeder Linie $A$ gleich viele Punkte $t$ liegen müssen (weil jede Linie $A$ in jede andere transformirbar ist), so ist, wenn $x$ die Anzahl dieser Punkte ist, $x . 21 : 3 = 28$, also $x = 4$. Ebenso schliesst man:

\gesperrt{Durch jeden Punkt} $a$ \gesperrt{gehen vier Linien} $T$\gesperrt{.}

Keiner der Punkte $t$ liegt auf einer Linie $T$.
Bezeichnen wir die $\nu$-zähligen Pole mit $P_\nu$, so haben wir also folgende Systeme von Polen:
\[
\begin{aligned}
&21 \text{ achtzählige Pole } P_8, \\
&28 \text{ sechszählige} \;\text{\textquotedbl}\; P_6, \\
&42 \text{ vierzählige} \;\text{\textquotedbl}\; P_4, \\
&56 \text{ dreizählige} \;\text{\textquotedbl}\; P_3, \\
&24 \text{ siebenzählige} \;\text{\textquotedbl}\; P_7, \\
&\text{unendlich viele zweizählige Pole } P_2.
\end{aligned}
\]

Die $P_2$ sind alle von $P_4, P_6, P_8$ verschiedene Punkte der Hauptaxen; von den Punkten $P_4$ liegen je zwei auf einer Linie $A$; von den Punkten $P_3$ liegen je zwei auf einer Linie $T$; die Punkte $P_7$ liegen nicht auf den Linien $A$ (dass sie auch nicht auf $T$ liegen, wird sich später ergeben).

Das System der Axen ist genau entsprechend dem System der Pole. Jede gerade Linie, die durch einen achtzähligen Pol geht, ist eine Axe (darunter die Hauptaxen, vier Linien $T$ und zwei Verbindungslinien des $P_8$ mit $P_4$). Demnach könnte
%CONT>>
